\documentclass[11pt]{article}

% --- Packages ---
\usepackage[utf8]{inputenc}
\usepackage[T1]{fontenc}
\usepackage{amsmath,amssymb,amsfonts}
\usepackage{graphicx}
\usepackage{booktabs}
\usepackage{hyperref}
\usepackage{url}
\usepackage{natbib}
\usepackage{xcolor}
\usepackage{microtype}
\usepackage{geometry}
\geometry{margin=1in}

\hypersetup{
    colorlinks=true,
    linkcolor=blue,
    citecolor=blue,
    urlcolor=blue
}

% --- Title ---
\title{Forced Single-Token Scoring for Multiple-Choice LLM Evaluation:\\
Why Extraction Strategy Dominates Prompt Engineering}

\author{
  Ari Choudhary \\
  UC Berkeley \\
  \texttt{ari@berkeley.edu}
}

\date{}

\begin{document}
\maketitle

% ===========================================================================
\begin{abstract}
We investigate inference-time optimization for HellaSwag commonsense reasoning using a locally-served 7B parameter model (qwen2.5:7b via Ollama). Our central finding is negative: standard prompt engineering techniques---chain-of-thought, self-consistency, few-shot learning---\emph{decrease} accuracy when applied to a free-text generation evaluation pipeline, with CoT causing a catastrophic 46.7 percentage point drop. We trace this failure to answer extraction brittleness, not model capability, and propose \textbf{forced single-token scoring}---constraining output to a single token after a completion-style prompt---as a simple fix that recovers +5.1\,pp immediately. Stacking TF-IDF few-shot selection and self-consistency majority voting on top of the corrected scoring yields a total improvement of +8.5\,pp (66.7\%~$\to$~75.2\%, $N{=}500$, McNemar $p{=}0.041$). We also evaluate choice-order shuffling (inspired by MedPrompt/SCOPE) as a position-bias mitigation strategy, finding it ineffective for this model (+0.6\,pp, $p{=}0.82$), suggesting qwen2.5:7b exhibits minimal positional preference. Our results align with recent findings by \citet{molfese2025} that MCQ evaluation strategy contributes more variance than prompt optimization, and that published prompt engineering gains may partially reflect extraction improvements rather than reasoning improvements.
\end{abstract}

% ===========================================================================
\section{Introduction}

The standard pipeline for evaluating LLMs on multiple-choice benchmarks involves three stages: (1)~format a prompt containing the question and answer choices, (2)~generate a response from the model, and (3)~extract the selected answer from the response. Most prompt engineering research focuses on stage~1 (better prompts) and stage~2 (better decoding), treating stage~3 as trivially solved by regex.

We show that stage~3---answer extraction---is the dominant source of error in generation-based MCQ evaluation, and that fixing it yields larger accuracy gains than any prompt optimization technique we tested. This echoes the concurrent finding by \citet{molfese2025} that ``traditional evaluation strategies often underestimate LLM capabilities'' due to extraction inconsistencies.

Our contributions:
\begin{itemize}
    \item We document a failure mode where every standard prompt optimization (CoT, few-shot, self-consistency) \emph{decreases} accuracy due to extraction brittleness.
    \item We propose forced single-token scoring as a simple, zero-cost fix.
    \item We demonstrate that TF-IDF few-shot selection and self-consistency stack on top of corrected scoring for a total +8.5\,pp improvement.
    \item We evaluate choice shuffling (MedPrompt/SCOPE-style) and present a negative result: it does not help when the model lacks strong position bias.
    \item We provide a complete, reproducible evaluation pipeline running on consumer hardware (Apple M3 Pro, 18\,GB RAM).
\end{itemize}

% ===========================================================================
\section{Related Work}

\paragraph{MCQ evaluation inconsistencies.}
\citet{molfese2025} systematically analyze MCQA evaluation strategies across three extraction methods (regex, logprobs, and LLM-based xFinder), four prompt settings (zero-shot, zero-shot CoT, zero-shot constrained, and few-shot), and eight models (1B--8B parameters) on MMLU-Redux (5{,}700 questions), OpenBookQA, and ARC-Challenge. Their central finding is that extraction method choice alone can shift reported accuracy by 3--6\,pp on the same model: on MMLU-Redux, regex yields 59.7\%, logprobs 60.5\%, and xFinder 62.3\% in the zero-shot setting. They also identify a fundamental trade-off: STEM subjects benefit from unconstrained prompts (+2.5\,pp with xFinder in ZS vs.\ ZS-Constrained), but unconstrained output raises regex miss rates above 40\%, and even the best LLM-based extractor (xFinder-Llama, 95.3\% Cohen's $\kappa$ with human raters) fails to detect inconsistent reasoning in adversarial examples at rates below 2\%.

Critically, \citeauthor{molfese2025} frame this as an unresolved trade-off---constrain the format and lose reasoning quality, or allow free generation and lose extraction reliability. \textbf{Our work resolves this trade-off.} Forced single-token scoring eliminates extraction error entirely (100\% parse rate) without sacrificing the model's ranking ability, because the next-token probability distribution $P(A \mid \text{prompt})$ already encodes the model's reasoning over the options. We do not need the model to \emph{articulate} its reasoning to \emph{apply} it. This is the key departure: where \citeauthor{molfese2025} conclude that better extractors are needed, we argue that the extraction problem can be sidestepped altogether by reading the model's beliefs directly from its output distribution.

Additionally, \citeauthor{molfese2025} find that few-shot prompting produces stable performance across all extraction strategies (converging to 63--64\% on MMLU-Redux), which aligns with our result that few-shot selection stacks cleanly on top of forced single-token scoring (+3.4\,pp). Their finding that response length beyond ${\sim}1{,}000$ characters provides only +0.6\% accuracy further supports our approach: the reasoning tokens are mostly waste from an evaluation standpoint.

Relatedly, \citet{molfese2025b} propose consistency evaluation with altered answer choices as a more reliable scoring methodology.

\paragraph{Multiple-choice evaluation.}
The lm-evaluation-harness \citep{gao2023harness} scores MCQ tasks using per-token log-likelihood: for each candidate continuation, compute $P(\text{continuation} \mid \text{context})$ and select the highest. This bypasses generation entirely. Our forced single-token approach approximates this through the generation interface when log-likelihoods are unavailable (as with many API-served models).

\paragraph{Chain-of-thought prompting.}
\citet{wei2022} showed that eliciting step-by-step reasoning improves performance on arithmetic and commonsense tasks. However, their evaluation used tasks with unambiguous numeric or short-form answers. We find CoT actively harms MCQ evaluation when the extraction pipeline cannot reliably parse the answer from verbose reasoning output.

\paragraph{Self-consistency.}
\citet{wang2023} proposed sampling multiple chain-of-thought paths and taking the majority vote. This assumes each sample produces a parseable answer. When extraction fails, self-consistency amplifies the failure mode---voting over 5 mis-extracted answers is worse than 1 mis-extracted answer.

\paragraph{Few-shot selection.}
\citet{liu2022} showed that selecting few-shot examples by semantic similarity outperforms random selection. We implement a lightweight TF-IDF variant that achieves similar benefits without requiring a neural embedding model, making it practical for local evaluation pipelines.

\paragraph{Constrained decoding.}
\citet{tam2024} study the impact of format restrictions on LLM performance and find that constrained generation can introduce artifacts. Our \texttt{num\_predict=1} approach is a minimal form of constrained decoding---rather than implementing logit masking, we simply limit generation length to 1 token after a completion-style prompt suffix, avoiding the complexity and potential artifacts of vocabulary restriction.

\paragraph{Position bias and choice shuffling.}
LLMs exhibit position bias in MCQ, disproportionately favoring certain answer slots regardless of content \citep{scope2025}. MedPrompt addresses this with choice-shuffling ensembles: shuffle answer order across $K$ runs, prompt with each variant, majority-vote. SCOPE formalizes this with an Inverse-Positioning module that estimates position preferences via null prompts. We evaluate vanilla choice shuffling ($K{=}5$) and find it ineffective for qwen2.5:7b, suggesting position bias severity is model-dependent.

\paragraph{Test-time compute scaling.}
\citet{snell2024} show that scaling inference-time compute (via majority voting, best-of-$N$, tree search) can be more effective than scaling model parameters. \citet{brown2024} find the optimal strategy varies by compute budget: shortest responses for low compute, beam search for medium, majority voting for high. Our self-consistency approach ($k{=}3$ majority voting) is a lightweight instance of test-time compute scaling, adding 3.5$\times$ latency for +3.4\,pp accuracy.

\paragraph{HyDE.}
\citet{gao2023hyde} proposed Hypothetical Document Embeddings---generating a hypothetical answer, embedding it, and using similarity search against real candidates. We considered this for HellaSwag but found that forced single-token scoring more directly accesses model beliefs with lower latency and no embedding model dependency.

% ===========================================================================
\section{Method}

\subsection{Problem Setup}

HellaSwag \citep{zellers2019} presents a context (e.g., ``A man picks up a knife in the kitchen'') and four candidate continuations labeled A--D. The task is to select the most plausible continuation. We evaluate using qwen2.5:7b (Q4\_K\_M quantization, 4.7\,GB) served locally via Ollama.

\subsection{Baseline: Free-Text Generation}

The naive approach prompts the model to choose:

\begin{verbatim}
Context: {context}
Which continuation is most natural?
A) {ending_0}
B) {ending_1}
C) {ending_2}
D) {ending_3}
Answer:
\end{verbatim}

The model generates free text (e.g., ``The answer is B because\ldots''). We extract the answer via regex, searching for the last standalone A/B/C/D character.

\textbf{Result:} 66.7\% accuracy ($N{=}20$). Extraction fails on ${\sim}15\%$ of examples where the model mentions multiple letters in its reasoning.

\subsection{Failed Optimizations on Free-Text}

\begin{table}[h]
\centering
\begin{tabular}{@{}llrl@{}}
\toprule
Technique & Accuracy & Delta & Failure Mode \\
\midrule
Baseline & 66.7\% & --- & --- \\
+ Chain-of-thought & 20.0\% & $-$46.7\,pp & Multi-paragraph output mentions all 4 letters \\
+ Self-consistency ($k{=}5$) & 30.0\% & $-$36.7\,pp & Majority vote over 5 mis-extractions \\
+ Few-shot ($k{=}3$) & 55.0\% & $-$11.7\,pp & Model hedges and discusses options \\
+ Prompt ensemble & 50.0\% & $-$16.7\,pp & Varying verbosity levels \\
\bottomrule
\end{tabular}
\caption{Every prompt optimization \emph{decreased} accuracy in the free-text generation setting. The prompt engineering was working---the extraction was breaking.}
\label{tab:failed}
\end{table}

\paragraph{Why self-consistency makes accuracy \emph{worse}, not just no-better.}
The drop from 66.7\% (baseline) to 30.0\% (self-consistency) is counterintuitive---majority voting should, in theory, only help. The failure mechanism has three parts:

\begin{enumerate}
    \item \textbf{Self-consistency inherently uses CoT + elevated temperature.} \citeauthor{wang2023}'s method samples diverse \emph{reasoning paths}, which requires temperature $> 0$ and chain-of-thought prompting. Both increase output length and verbosity, which is exactly what breaks regex extraction. Each of the $k{=}5$ samples suffers from the same class of extraction failure as standalone CoT (20.0\%).

    \item \textbf{Extraction errors are correlated, not independent.} Majority voting's theoretical guarantee assumes each sample independently produces a correct answer with probability $p$, so voting pushes aggregate accuracy above $p$. But extraction errors are \emph{systematic}: when the model discusses all four options in order (A, then B, then C, then D), the regex consistently latches onto the same wrong letter across samples---typically whichever letter appears last in the reasoning or first in a phrase like ``Option A is unlikely.'' Because the same prompt structure drives all $k$ samples toward similar verbose patterns, the extraction errors are positively correlated.

    \item \textbf{Voting amplifies correlated errors.} If 3 out of 5 samples all mis-extract to the same wrong letter (say ``A''), the majority vote locks in ``A'' with high confidence. A single greedy sample at temperature$\,{=}\,$0 at least has a ${\sim}85\%$ chance of producing concise-enough output for successful extraction; with self-consistency, the vote is dominated by the correlated mis-extractions. The result (30.0\%) is above CoT alone (20.0\%) because voting occasionally recovers when 2--3 samples happen to extract correctly, but far below baseline because the per-sample extraction success rate (${\sim}30$--$40\%$) is too low for majority voting to reliably converge on the right answer.
\end{enumerate}

In short: self-consistency assumes a reliable answer-extraction channel. When that channel is noisy and biased, voting doesn't average out the noise---it amplifies the bias.

\subsection{Forced Single-Token Scoring}

We modify the prompt to end with a completion-style suffix and constrain output to exactly one token:

\begin{verbatim}
Context: {context}
Which continuation is most natural?
A) {ending_0}  B) {ending_1}  C) {ending_2}  D) {ending_3}
The answer is
\end{verbatim}

With \texttt{num\_predict=1}, the model outputs exactly one token (typically ``A'', ``B'', ``C'', or ``D''). Extraction is trivial: read the first alphabetic character.

\paragraph{Why this works.} The probability distribution $P(\text{next\_token} \mid \text{prompt})$ already encodes the model's reasoning. By preventing generation of reasoning tokens, we:
\begin{enumerate}
    \item Eliminate extraction failures entirely.
    \item Force the model to commit to its highest-confidence answer.
    \item Reduce latency from ${\sim}4.6$\,s to ${\sim}0.5$\,s per example (1 token vs.\ ${\sim}100$ tokens).
\end{enumerate}

This is equivalent to constrained decoding with a vocabulary restricted to $\{A, B, C, D\}$, implemented without logit manipulation.

\subsection{Stacked Optimizations}

With extraction fixed, standard techniques work as intended:

\paragraph{TF-IDF few-shot selection ($k{=}3$).}
For each test example, we select the 3 most similar training examples by TF-IDF cosine similarity on the context text. We build a document-frequency index over 500 training examples, compute TF-IDF vectors (no sklearn---pure Python), and rank by cosine similarity. Selected examples are prepended in the same forced-choice format, demonstrating the expected single-token response pattern.

\paragraph{Self-consistency ($k{=}3$, temp$\,{=}\,$0.3).}
We generate 3 responses at temperature$\,{=}\,$0.3 with seeds $\{42, 49, 56\}$ and take the majority vote. With single-token output, each sample is guaranteed to produce a valid A/B/C/D answer, so voting is reliable.

\subsection{Choice Shuffling (MedPrompt-style)}

For each example, we generate $K{=}5$ random permutations of the answer order, prompt the model with each permuted version, map each prediction back to the original answer index, and majority-vote across the $K$ runs. This is designed to cancel position bias---if the model always favors ``A'', shuffling ensures ``A'' maps to different original answers across runs.

\subsection{Statistical Testing}

We use McNemar's test on paired predictions (baseline vs.\ optimized on the same 500 examples). McNemar's is appropriate because the samples are dependent---each example is evaluated under both conditions. We report the one-sided exact binomial test ($H_1$: optimized $>$ baseline) and Wilson score confidence intervals.

% ===========================================================================
\section{Experiments}

\subsection{Setup}

\begin{itemize}
    \item \textbf{Model:} qwen2.5:7b (Q4\_K\_M, 4.7\,GB) via Ollama
    \item \textbf{Hardware:} Apple M3 Pro, 18\,GB unified memory
    \item \textbf{Dataset:} HellaSwag validation split (500 examples) + training split (500 examples for few-shot pool)
    \item \textbf{Seeds:} Baseline: 42. Self-consistency: 42, 49, 56. Shuffle: 42 (deterministic permutation generator).
    \item \textbf{Metrics:} Accuracy, Wilson 95\% CI, bootstrap 95\% CI (2{,}000 resamples), McNemar $p$-value, latency
\end{itemize}

\subsection{Main Results}

\begin{table}[h]
\centering
\begin{tabular}{@{}lcccc@{}}
\toprule
Method & Accuracy & 95\% CI (Wilson) & Avg Latency & $p$ vs baseline \\
\midrule
Free-text baseline & 66.7\% ($N{=}20$) & [48.8\%, 80.8\%] & 4.59\,s & --- \\
Single-token baseline & 71.8\% ($N{=}500$) & [67.7\%, 75.6\%] & 0.46\,s & --- \\
Single-token + shuffle ($K{=}5$) & 72.4\% ($N{=}500$) & [68.3\%, 76.1\%] & 1.82\,s & 0.82 \\
\textbf{Single-token + few-shot + SC} & \textbf{75.2\%} ($N{=}500$) & \textbf{[71.2\%, 78.8\%]} & \textbf{1.63\,s} & \textbf{0.041} \\
\bottomrule
\end{tabular}
\caption{Main results. The total improvement decomposes as: scoring method fix (free-text $\to$ single-token): \textbf{+5.1\,pp}; few-shot + self-consistency: \textbf{+3.4\,pp}; total: \textbf{+8.5\,pp}.}
\label{tab:main}
\end{table}

\subsection{Choice Shuffling: A Negative Result}

Choice shuffling ($K{=}5$ permutations, majority vote) yielded only +0.6\,pp over single-token baseline (72.4\% vs 71.8\%), with McNemar $p{=}0.82$---not significant.

\begin{table}[h]
\centering
\begin{tabular}{@{}lr@{}}
\toprule
Category & Count \\
\midrule
Both correct & 322 \\
Shuffle better & 40 \\
Baseline better & 37 \\
Both wrong & 101 \\
\bottomrule
\end{tabular}
\caption{Discordant pair analysis for choice shuffling vs.\ baseline.}
\label{tab:shuffle}
\end{table}

The near-equal split of discordant pairs (40 vs 37) indicates that shuffling adds noise rather than correcting bias. This suggests \textbf{qwen2.5:7b does not exhibit strong position bias} on HellaSwag---the model's errors come from genuine reasoning limitations, not option-order sensitivity. This contrasts with larger models (GPT-4, Claude) where position bias is well-documented and shuffling helps significantly (MedPrompt reports +2--4\,pp on medical MCQ).

\subsection{Discordant Pair Analysis (Few-Shot + SC vs Baseline)}

McNemar's test on the 500 paired examples:

\begin{table}[h]
\centering
\begin{tabular}{@{}lr@{}}
\toprule
Category & Count \\
\midrule
Both correct & 325 \\
Fixed (baseline wrong $\to$ optimized correct) & 51 \\
Regressed (baseline correct $\to$ optimized wrong) & 34 \\
Both wrong & 90 \\
\bottomrule
\end{tabular}
\caption{Discordant pair analysis for few-shot + self-consistency vs.\ baseline.}
\label{tab:discordant}
\end{table}

Net: +17 correct. One-sided exact binomial $p = 0.041$.

\subsection{Error Analysis}

\paragraph{Fixed examples (51).} The most common pattern is the baseline predicting ``A'' as a default (the model's prior for the first option) while the optimized version, primed by similar few-shot examples, correctly distributes probability mass. Example~\#8: ``A man holding a pocket knife while sitting on rocks\ldots''---baseline defaults to A, optimized correctly picks B.

\paragraph{Regressed examples (34).} Few-shot examples occasionally prime the model toward a wrong pattern. Example~\#0: ``A man sitting on a roof\ldots''---the few-shot examples involve indoor activities, biasing toward an indoor continuation (B) when the correct answer is outdoor (D).

\paragraph{Both wrong (90).} These are genuinely ambiguous or require world knowledge the 7B model lacks. Example~\#6: ``A man playing a harmonica\ldots''---both predict A, gold is C. The correct continuation involves a specific musical technique the model hasn't seen enough training data for.

\subsection{Latency Analysis}

\begin{table}[h]
\centering
\begin{tabular}{@{}lccc@{}}
\toprule
Configuration & Requests/example & Avg Latency & Throughput \\
\midrule
Single-token baseline & 1 & 0.46\,s & 2.17 ex/s \\
+ Shuffle ($K{=}5$) & 5 & 1.82\,s & 0.55 ex/s \\
+ Few-shot + SC ($k{=}3$) & 3 & 1.63\,s & 0.61 ex/s \\
\bottomrule
\end{tabular}
\caption{Latency comparison. Few-shot + SC achieves higher accuracy than shuffle at lower latency.}
\label{tab:latency}
\end{table}

\subsection{Additional Benchmarks}

\begin{table}[h]
\centering
\begin{tabular}{@{}lcc@{}}
\toprule
Benchmark & Accuracy & $N$ \\
\midrule
MMLU (6 subjects) & 65.3\% & 300 \\
Custom (prog\_lang\_id) & 100.0\% & 18 \\
\bottomrule
\end{tabular}
\caption{Additional benchmark results using forced single-token scoring.}
\label{tab:additional}
\end{table}

% ===========================================================================
\section{Alternative Approaches Considered}

\subsection{Log-Probability Scoring via lm-eval}

The standard approach in lm-eval computes $P(\text{continuation} \mid \text{context})$ using per-token log-probabilities. We implemented this using Ollama's \texttt{/v1/chat/completions} endpoint with \texttt{logprobs=true}, scoring 5 tokens per continuation via token-by-token probing with an assistant-prefix message. However, the chat template wrapping introduces noise---the assistant-prefix approach doesn't perfectly replicate raw completion scoring. Result: 38\% accuracy on 50 examples, well below our forced single-token approach (72\%). This aligns with the finding that serving infrastructure gaps (logprob access) are a primary barrier to reproducible local evaluation.

\subsection{HyDE (Hypothetical Document Embeddings)}

Generate a hypothetical continuation, embed it, embed each real choice, pick the closest by cosine similarity. We didn't pursue this because: (1)~qwen2.5:7b isn't an embedding model---a dedicated model like \texttt{nomic-embed-text} would be needed, (2)~6 API calls per example vs.\ our 3, (3)~forced single-token more directly accesses model beliefs.

\subsection{Logprob-Based Confidence Filtering}

Use the logprob of the generated token as a confidence score. Abstain on low-confidence examples and evaluate only high-confidence ones. This trades coverage for precision and could be valuable in deployment but doesn't directly improve accuracy on the full test set.

\subsection{Shortest Majority Vote}

Recent work on test-time compute scaling \citep{snell2024} shows that combining parallel scaling strategies with response-length characteristics can improve scalability. ``Shortest Majority Vote'' selects the most common answer among the shortest CoT responses, exploiting the observation that correct CoT solutions tend to be more concise. This is incompatible with our single-token approach (all responses are 1 token) but would be relevant for free-text evaluation pipelines.

\subsection{Compute-Optimal Inference Strategy Selection}

\citet{brown2024} find that the optimal test-time strategy varies by compute budget: greedy for minimal compute, beam search for medium, majority voting for high. Our results are consistent: at 1 request/example (minimal compute), greedy single-token performs well (71.8\%); at 3 requests/example (moderate compute), majority voting adds +3.4\,pp. The diminishing returns suggest that for this model and task, further scaling ($k{=}5$, $k{=}10$) would yield marginal gains.

% ===========================================================================
\section{Calls to Action}

\begin{enumerate}
    \item \textbf{Benchmark authors should specify extraction methodology.} Published accuracy numbers are only meaningful if the extraction pipeline is specified. \citet{molfese2025} demonstrate that different extraction methods produce accuracy differences of 10+ pp on the same model. A 10-point accuracy difference can come entirely from how you parse the model's output, not from model capability.

    \item \textbf{Report extraction failure rates alongside accuracy.} If 15\% of your model's outputs can't be parsed into a valid answer, that's a measurement problem, not a model problem. Extraction failure rate should be a standard metric in MCQ evaluation.

    \item \textbf{Default to constrained decoding for MCQ.} Free-text generation + regex is a legacy pattern from when models were instruction-tuned to output clean answers. Modern models with CoT capabilities will reason at length unless constrained. \texttt{num\_predict=1} or logit masking should be the default for MCQ evaluation.

    \item \textbf{Separate ``reasoning improvement'' from ``extraction improvement'' in ablations.} If your CoT optimization improves accuracy by 5\%, verify that it's not simply because the model now outputs the letter on the last line instead of buried in a paragraph. Run the same ablation with constrained decoding to isolate the reasoning contribution.

    \item \textbf{Test position bias before applying choice shuffling.} Our negative result with shuffling shows that position bias is model-dependent. SCOPE's null-prompt probing approach can cheaply estimate bias severity before committing to the 5$\times$ latency cost of shuffling.

    \item \textbf{Standardize logprob APIs across serving frameworks.} The gap between API-based and locally-served evaluation is primarily in logprob access, not model quality. Ollama's lack of prompt-token logprobs forced us to develop alternative scoring strategies. A standardized logprob interface (like OpenAI's \texttt{logprobs} parameter) across Ollama, vLLM, and TGI would improve reproducibility.
\end{enumerate}

% ===========================================================================
\section{Scope and Limitations: When CoT Genuinely Helps}

Our results should not be read as a general indictment of chain-of-thought prompting. The claim is narrower: \textbf{for multiple-choice ranking tasks like HellaSwag, CoT adds extraction noise without adding reasoning signal}, because the model's belief about which continuation is most plausible is already captured in its single-token output distribution $P(A \mid \text{prompt})$.

This does \emph{not} generalize to tasks requiring multi-step compositional reasoning. Consider a complex algebraic problem: if you first prompt the model to enumerate the properties of the relevant structure (e.g., ``list the properties of a nilpotent Lie algebra''), those properties are serialized into the context window and become explicit conditioning for subsequent generation. The model no longer needs to implicitly recall these properties through its weights---it reads them back from its own output. Each intermediate step narrows the distribution for the next, functioning as an external scratchpad that compensates for the transformer's bounded per-pass computation depth.

This scratchpad effect is well-established \citep{wei2022} and is the genuine mechanism behind CoT's success on arithmetic, symbolic reasoning, and multi-hop QA. The key distinction is between:

\begin{itemize}
    \item \textbf{Recognition/ranking tasks} (HellaSwag, most MCQ benchmarks): The model compares pre-stated options. No intermediate computation needs to be serialized. The output distribution already encodes the ranking. CoT is overhead.
    \item \textbf{Generation/reasoning tasks} (math, logic, planning): The answer depends on intermediate results that exceed what the model can compute in a single forward pass. CoT is essential---it externalizes working memory into the context window, and forced single-token scoring would be catastrophic.
\end{itemize}

Our forced single-token approach is a measurement technique for the first category. It should not be applied to the second. The broader lesson is that evaluation methodology must be matched to task structure: constrained decoding for recognition, free generation for reasoning.

% ===========================================================================
\section{Conclusions and Future Work}

We demonstrated that answer extraction methodology is the dominant variable in multiple-choice LLM evaluation, contributing more accuracy variance than any prompt engineering technique we tested. Forced single-token scoring provides a simple, zero-cost fix that eliminates extraction failures and enables standard prompt optimizations to work as intended. Choice shuffling, while theoretically motivated, does not help for models without strong position bias.

\paragraph{Key numbers:}
\begin{itemize}
    \item Free-text $\to$ single-token: \textbf{+5.1\,pp} (extraction fix)
    \item + Few-shot + self-consistency: \textbf{+3.4\,pp} (prompt optimization)
    \item Choice shuffling: \textbf{+0.6\,pp} (not significant, $p{=}0.82$)
    \item Total: \textbf{+8.5\,pp} (66.7\% $\to$ 75.2\%, $p = 0.041$)
\end{itemize}

\subsection{Future Work}

\begin{enumerate}
    \item \textbf{SCOPE-style structured shuffling.} Our shuffling used random permutations. SCOPE's Inverse-Positioning module uses null prompts to estimate position preferences and Semantic-Spread to optimize option placement. This structured approach may outperform random shuffling even for models with weak bias.

    \item \textbf{Adaptive test-time compute.} Following \citet{snell2024}, use the logprob of the single generated token as a confidence signal. Apply self-consistency only on low-confidence examples (logprob $<$ threshold), and accept greedy answers on high-confidence ones. This could achieve the accuracy of $k{=}3$ at the average latency of $k{\sim}1.5$.

    \item \textbf{Logprob-based scoring with vLLM.} Replace Ollama with vLLM to get native prompt-token log-probabilities. This would enable proper perplexity-based scoring and likely outperform our generation-based approach. vLLM's PagedAttention also enables higher throughput for batch evaluation.

    \item \textbf{Cross-model position bias profiling.} Run SCOPE's null-prompt probing across model families (Llama~3, Mistral, Phi, Qwen) and sizes (1B--70B) to build a taxonomy of which models benefit from shuffling and which don't.

    \item \textbf{Combined shuffling + few-shot.} We tested shuffling and few-shot independently. Combining them (shuffle the few-shot examples' answer order too, matching the test question's shuffled order) could provide both position-bias correction and semantic priming simultaneously.

    \item \textbf{Weighted majority voting with PRM.} Recent work on optimal aggregation of LLM and Process Reward Model (PRM) signals shows that a weighted majority vote combining generation confidence with step-level verification outperforms uniform voting. Integrating a lightweight PRM into the self-consistency pipeline could improve accuracy without additional model calls.

    \item \textbf{Embedding-based scoring (HyDE).} Using a dedicated embedding model (e.g., \texttt{nomic-embed-text} via Ollama) alongside the generation model for hypothesis-then-match scoring. Particularly promising for tasks where the correct answer is semantically distinct from distractors.
\end{enumerate}

% ===========================================================================
\bibliographystyle{plainnat}

\begin{thebibliography}{99}

\bibitem[Brown et~al.(2024)]{brown2024}
Brown, B. et~al. (2024).
\newblock Inference Scaling Laws: An Empirical Analysis of Compute-Optimal Inference for Problem-Solving with Language Models.
\newblock \emph{ICLR 2025}. \href{https://arxiv.org/abs/2408.00724}{arXiv:2408.00724}.

\bibitem[Gao et~al.(2023a)]{gao2023harness}
Gao, L. et~al. (2023a).
\newblock A framework for few-shot language model evaluation.
\newblock lm-evaluation-harness. \url{https://github.com/EleutherAI/lm-evaluation-harness}.

\bibitem[Gao et~al.(2023b)]{gao2023hyde}
Gao, L. et~al. (2023b).
\newblock Precise Zero-Shot Dense Retrieval without Relevance Labels (HyDE).
\newblock \emph{ACL 2023}.

\bibitem[Liu et~al.(2022)]{liu2022}
Liu, J. et~al. (2022).
\newblock What Makes Good In-Context Examples for GPT-3?
\newblock \emph{DeeLIO Workshop, ACL 2022}.

\bibitem[Molfese et~al.(2025a)]{molfese2025}
Molfese, F.~M. et~al. (2025a).
\newblock Right Answer, Wrong Score: Uncovering the Inconsistencies of LLM Evaluation in Multiple-Choice Question Answering.
\newblock \emph{ACL 2025 Findings}. \href{https://arxiv.org/abs/2503.14996}{arXiv:2503.14996}.

\bibitem[Molfese et~al.(2025b)]{molfese2025b}
Molfese, F.~M. et~al. (2025b).
\newblock Improving Score Reliability of Multiple Choice Benchmarks with Consistency Evaluation and Altered Answer Choices.
\newblock \href{https://arxiv.org/abs/2511.21860}{arXiv:2511.21860}.

\bibitem[SCOPE(2025)]{scope2025}
SCOPE (2025).
\newblock Stochastic and Counterbiased Option Placement for Evaluating Large Language Models.
\newblock \href{https://arxiv.org/abs/2507.18182}{arXiv:2507.18182}.

\bibitem[Snell et~al.(2024)]{snell2024}
Snell, C. et~al. (2024).
\newblock Scaling LLM Test-Time Compute Optimally can be More Effective than Scaling Model Parameters.
\newblock \href{https://arxiv.org/abs/2408.03314}{arXiv:2408.03314}.

\bibitem[Tam et~al.(2024)]{tam2024}
Tam, Z.~R. et~al. (2024).
\newblock Let Me Speak Freely? A Study on the Impact of Format Restrictions on LLM Performance.

\bibitem[Wang et~al.(2023)]{wang2023}
Wang, X. et~al. (2023).
\newblock Self-Consistency Improves Chain of Thought Reasoning in Language Models.
\newblock \emph{ICLR 2023}.

\bibitem[Wei et~al.(2022)]{wei2022}
Wei, J. et~al. (2022).
\newblock Chain-of-Thought Prompting Elicits Reasoning in Large Language Models.
\newblock \emph{NeurIPS 2022}.

\bibitem[Zellers et~al.(2019)]{zellers2019}
Zellers, R. et~al. (2019).
\newblock HellaSwag: Can a Machine Really Finish Your Sentence?
\newblock \emph{ACL 2019}.

\end{thebibliography}

\vspace{1em}
\noindent\rule{\textwidth}{0.4pt}
\footnotesize{All experiments conducted on Apple M3 Pro (18\,GB) with qwen2.5:7b (Q4\_K\_M) via Ollama. Code and results available at \url{https://github.com/arihantchoudhary/mercor}.}

\end{document}
